\section{问题重述}

\subsection{问题背景}

我国社区养老服务体系正在由机构集中供给逐步转向社区嵌入式、家门口可及和医养康养结合的综合服务模式。对成熟街道而言，养老服务站不仅要解决老人能否到达的问题，还要解决服务能力是否足够、服务价格是否可承受、机构经营是否可持续等问题。某街道下辖 10 个连片小区，政府计划在小区内部建设若干嵌入式社区养老服务站，为老人提供助餐、日间照料、上门护理、康复理疗、助浴和紧急救助等服务。题目提供了小区基础数据、服务需求数据、站点建设运营成本、小区间距离矩阵和满意度评分规则。本文需要基于这些数据建立预测、选址、定价和鲁棒性评价模型。

\subsection{问题提出}

本文需要解决四个问题。第一，基于当前各小区老人数量、状态转移概率、死亡率和新增老人比例，预测未来 5 年各小区三类老人数量，并计算第 5 年末理论月服务需求和消费约束后的有效需求。第二，在总建设预算不超过 120 万元、有效服务半径不超过 1000 米的约束下，确定服务站建设数量、位置和规模，使覆盖率和满意度尽可能高，并计算站点利润和小区满意度。第三，在问题二确定的选址方案基础上，考虑政府补贴和利润率约束，确定各项服务价格，使老人满意度最大化。第四，改变老人增长率、失能转移概率、固定管理成本和建设预算，重新求解问题二和问题三，比较方案变化并评价模型鲁棒性。
